% engineering/part_intro.tex

\chapter*{Engineering}
\phantomsection
\addcontentsline{toc}{chapter}{Engineering}

% ============================================================
\section*{What do we already know?}
% ============================================================

\begin{remark}
The previous parts established the mathematical, geometric,
operational, realization-aware, and optimization foundations of AIM.
\end{remark}

\begin{remark}
Railway alignments can now be represented, evaluated, analyzed, and
improved through structured operator-based methods.
\end{remark}

\begin{remark}
Optimization provides mechanisms for generating improved solutions
according to specified objectives.
\end{remark}

\begin{remark}
However, optimization alone does not determine whether a solution is
acceptable from an engineering perspective.
\end{remark}

% ============================================================
\section*{What is still missing?}
% ============================================================

\begin{remark}
Railway engineering operates within a framework of rules,
recommendations, standards, regulations, and operational constraints.
\end{remark}

\begin{remark}
Not every mathematically feasible alignment is an admissible railway
alignment.
\end{remark}

\begin{remark}
Track geometry must satisfy engineering limits.

Operational states must remain safe.

Construction solutions must remain practical.

Design alternatives must comply with applicable standards.
\end{remark}

\begin{remark}
The missing element is therefore the engineering constraint layer that
connects mathematical models with railway practice.
\end{remark}

% ============================================================
\section*{What comes next?}
% ============================================================

\begin{remark}
This part introduces engineering constraints as explicit components of
the AIM framework.
\end{remark}

\begin{remark}
Rather than treating standards as external documents consulted after
design, AIM interprets engineering rules as structured constraints that
can participate directly in analysis and optimization.
\end{remark}

\begin{remark}
The chapters that follow investigate:
\begin{itemize}
\item standards and regulations,
\item admissibility constraints,
\item railway-specific engineering rules,
\item and formal representations of operational requirements.
\end{itemize}
\end{remark}

\begin{remark}
The objective is to transform railway standards from passive reference
material into active components of railway information modelling.
\end{remark}

% ============================================================
\section*{Relation to the AIM Roadmap}
% ============================================================

\begin{remark}
Within the AIM roadmap, the present part corresponds to the transition
from optimization toward engineering admissibility.
\end{remark}

\begin{center}
\begin{tikzpicture}[
  node distance=1.4cm,
  box/.style={
    draw,
    rounded corners,
    minimum width=5.0cm,
    minimum height=0.9cm,
    align=center
  }
]

\node[box] (optimization)
{Optimization};

\node[box, below=of optimization] (constraints)
{Engineering Constraints};

\node[box, below=of constraints] (rules)
{Railway Rules};

\draw[->, thick] (optimization) -- (constraints);
\draw[->, thick] (constraints) -- (rules);

\end{tikzpicture}
\end{center}

\begin{remark}
The progression
\[
\text{Optimization}
\rightarrow
\text{Engineering Constraints}
\rightarrow
\text{Railway Rules}
\]
marks the transition from finding improved solutions toward ensuring
that those solutions remain admissible within real railway systems.
\end{remark}

\begin{remark}
Engineering therefore acts as the bridge between mathematical
optimization and practical railway implementation.
\end{remark}

% ============================================================
\begin{principle}[Engineering Constraint Principle]
\label{princ:engineering-constraint-principle}
Within AIM, railway standards, regulations, and operational rules
should be represented as explicit engineering constraints that can
participate directly in modelling, analysis, and optimization.
\end{principle}

% ============================================================
\begin{summary}
Optimization explains how railway systems can be improved.

Engineering explains which improvements are admissible.

The present part therefore establishes the constraint layer that
connects mathematical optimization with practical railway design and
operation.
\end{summary}
